\scanpage{149}
résulte que l'on a un plus petit couple $(E_0,K_0)$ $(E_0\subset E,\ K_0\subset K)$
tel que $\mathfrak U$ provienne d'un idéal de $E_0\otimes K_0$.
On a vu que $E_0,K_0$ sont weiliens non réels. Si $E_0$ est fini sur $\mathbb Q$,
de degré $2d$, alors $K_0$ est fini sur $\mathbb Q$, de degré $\leq\binom{2d}{d}$
[car le ss-groupe des $g\in G$ qui invarient une partie à $d$ éléments
dans un ens. à $2d$ éléments est d'indice $\leq\binom{2d}{d}$].
\deleted{Si $K$ f $K$} Si $E$ fini sur $\mathbb Q$ de degré $2d$, alors
$\mathfrak A(\Sigma):=(E\otimes K)/\mathfrak U$ est fini sur $K$ de degré $d$,
et réciproqt.

Soient $E,K$ deux extensions de $\mathbb Q$, $t$ un bimodule sur $E,K$, i.e. un module
sur $E\otimes K$. On dit que $t$ est du type CM si $t$ est monogène, et son annulateur
est un $\mathfrak U$ de type CM. Les bimodules de type CM \inserted{sur $E,K$} sont
classifiés, à isom près, par les structures CM $\deleted{\mathfrak U}\Sigma$ sur le couple $E,K$.

Un tel $t$ définit \inserted{(si \deleted{E} $K$ fini sur $\mathbb Q$, i.e. $t$ fini sur $E$),}
un hom
\[
\psi_\Sigma:T_K\longrightarrow T_E
\]
dont l'action, sur les pts rationnels sur $\mathbb Q$ (pour fixer les idées) est donnée par
\[
\psi_\Sigma(x)=\operatorname{dét}_{t/E}(x_t)\,;
\]
Ou si on préfère, en termes de la fonction caractéristique de $\Sigma$
(qui est ouvert et fermé), notée \deleted{aussi $\Sigma$} $1_\Sigma$
\[
\psi_\Sigma(x)=N_{E\otimes K/E}(x1_\Sigma^{\deleted{\Sigma}})
\]
Si $K'$ est une extension finie de $K$, \inserted{et $E'$ ext. de $E$,}
donnant lieu à un $\Sigma'$ dans $E'\otimes K'$, on aura commutativité dans
\diagram{\begin{tikzcd}[column sep=large]
T_{K'} \arrow[r,"N_{K'/K}"] \arrow[rrr,bend right=24,"\psi_{\Sigma'}"'] &
T_K \arrow[r,no head,"\psi_\Sigma"] & T_E \arrow[r,no head,"i_{E'/E}"] & T_{E'}
\end{tikzcd}}
\margin{148}
